% Mini-Lang：以 256 字受限中文研究小型語言模型的泛化
% 編譯：xelatex paper.tex（執行兩次以更新交互參照）
\documentclass[11pt,a4paper]{article}

\usepackage[margin=25mm]{geometry}
\usepackage{fontspec}
\usepackage{amsmath,amssymb}
\usepackage{booktabs,tabularx,array}
\usepackage{float}
\usepackage{caption}
\usepackage{xcolor}
\usepackage{tikz}
\usetikzlibrary{arrows.meta,positioning}
\usepackage[hidelinks]{hyperref}

% ---- 中文字型：有 xeCJK 就用 xeCJK，否則用 fontspec + XeTeX 內建的中文斷行 ----
\IfFileExists{xeCJK.sty}{%
  \usepackage{xeCJK}
  \setCJKmainfont{Songti TC}[BoldFont={Songti TC Bold}, ItalicFont={Heiti TC}]
  \setCJKsansfont{Heiti TC}
  \setCJKmonofont{Heiti TC}
}{%
  \setmainfont{Songti TC}[BoldFont={Songti TC Bold}, ItalicFont={Heiti TC}]
  \setsansfont{Heiti TC}
  \XeTeXlinebreaklocale "zh"
  \XeTeXlinebreakskip = 0pt plus 1pt minus 0.1pt
}

\linespread{1.3}
\setlength{\parindent}{2em}
\setlength{\parskip}{0.3em}
\captionsetup{font=small,labelsep=quad}
\renewcommand{\figurename}{圖}
\renewcommand{\tablename}{表}
\renewcommand{\abstractname}{摘要}
\renewcommand{\refname}{參考文獻}
\renewcommand{\contentsname}{目錄}
\renewcommand{\appendixname}{附錄}
\newcolumntype{Y}{>{\raggedright\arraybackslash}X}
\newcommand{\tok}[1]{\allowbreak\texttt{<#1>}}

\title{\textbf{Mini-Lang：以 256 字受限中文研究小型語言模型的泛化}\\[0.6em]
\large Mini-Lang: Studying Generalization of Small Language Models\\ with a 256-Character Controlled Chinese}
\author{ccckmit\thanks{待補：姓名、單位與聯絡方式。}}
\date{}

\begin{document}
\maketitle

\begin{abstract}
大型語言模型的能力常被歸因於規模，但規模同時混雜了模型大小、資料量與資料多樣性三個因素。本文設計一個受限的中文語言 \textbf{Mini-Lang}：只用 256 個漢字、17 條固定語法規則，而且每個合法句子都是自然的標準中文。在這個「方向數得完」的小語言上，我們用約 4 萬到 640 萬參數的 Transformer 進行一系列可控實驗。主要發現如下：
(1) 只用模板資料時，小模型會死背：在 50 句的蒸餾語料上，訓練集全對，驗證集為 0\%；
(2) 縮小模型在每一種資料量下都更差：640 萬參數的模型搭配 16 倍資料，在新人物事實問答上達到 98.6\%，4 萬參數的模型則和「猜最常見答案」的基準相當；
(3) 正則化只在資料少時有益：在 0.5 倍資料時提升 27 個百分點，在 16 倍資料時下降 57 個百分點；
(4) 混入的語句只幫助它所代表的方向：加入同義問法使「換一種問法」的準確率從 68.5\% 升到 85.5\%；
(5) 以大型語言模型依 13 個方向產生約 1,800 句語料（以句數計約佔預訓練語料 10\%），沒看過的句子的每字 loss 從 6.74 降到 3.30，代價是部分問答準確率下降。
此外，我們讓兩個知識不同的 Mini-Lang 模型互相問答：所有 80 句話都只用到 256 字，但從對話中學到的知識只限於聽過的題目，並沒有推廣到同一人物的其他事實。

\medskip
\noindent\textbf{關鍵詞：}受限語言、小型語言模型、泛化、資料多樣性、縮放實驗、合成資料
\end{abstract}

% ============================================================
\section{緒論}

大型語言模型（LLM）展現出許多「湧現」能力 \cite{wei2022emergent}，但這些能力究竟來自模型大小、資料量，還是資料的多樣性，在大規模實驗中很難拆開。另一方面，小型語言模型常被認為「泛化能力不足」，可是這種說法很少在可控的條件下被檢驗。

本文的出發點是 Ogden 在 1930 年提出的 Basic English \cite{ogden1930basic}：只用 850 個英文詞，就能表達日常生活中大部分的意思。TinyStories \cite{eldan2023tinystories} 進一步顯示，只要把語言限制在幼童能懂的詞彙，並提供大量多樣的故事，100 萬到 3,000 萬參數的模型也能寫出連貫的英文。

我們把這個思路用在中文。首先設計 \textbf{Mini-Lang}，一個只用 256 個漢字的中文子集：256 字剛好可以用一個位元組編碼，詞表固定，所有方向都數得完。接著在 Mini-Lang 上建立一個可以完全控制的實驗環境：語料由程式產生，人物、事實、問法、句型、資料量都能任意調整，而且每個問題都有標準答案。我們用這個環境回答幾個問題：小模型為什麼泛化不足？是模型太小，還是資料不夠？把模型縮小能不能逼它學規則？混入額外的語句，能讓模型在哪些方向泛化得更好？兩個小模型能不能透過對話交換知識？

本文的貢獻如下：
\begin{enumerate}
  \item \textbf{Mini-Lang}：一個 256 字、17 條語法規則的受限中文，附有自動檢查工具，確保語料只用到字表內的字（第 \ref{sec:minilang} 節）。
  \item \textbf{一個完整的小型實驗平台}：包含字元級 Transformer、多語料訓練流程、可控的語料產生器，以及兩個模型互相交談的 mini-talk（第 \ref{sec:system} 節）。
  \item \textbf{一系列可重現的實驗}：縮放實驗（4 種模型大小 $\times$ 4 種資料量，另加正則化）、資料設計實驗（4 種語料設定 $\times$ 4 種泛化測試），以及對話學習實驗（第 \ref{sec:exp} 節）。
  \item \textbf{對泛化的實證分析}：區分分布內與分布外泛化，並指出「定向覆蓋」與「容量」的作用（第 \ref{sec:discussion} 節）。
\end{enumerate}

% ============================================================
\section{相關研究}

\paragraph{受限語言。}
Basic English \cite{ogden1930basic} 以 850 個詞構成英文的子集。TinyStories \cite{eldan2023tinystories} 以 GPT-4 產生只用簡單詞彙的故事，證明極小的模型也能學會流暢、連貫的語言。Mini-Lang 承襲這兩者：它和 Basic English 一樣是自然語言的子集，又像 TinyStories 一樣用大型語言模型產生多樣的語料。

\paragraph{縮放定律與湧現。}
Kaplan 等人 \cite{kaplan2020scaling} 發現 loss 隨參數量、資料量與計算量呈冪次下降，且大模型的樣本效率較高。Chinchilla \cite{hoffmann2022chinchilla} 建議資料量約為參數量的 20 倍。Wei 等人 \cite{wei2022emergent} 描述了隨規模突然出現的能力；Schaeffer 等人 \cite{schaeffer2023mirage} 則指出，這種「突然」有一部分來自全有全無的評分方式。

\paragraph{泛化與記憶。}
雙重下降 \cite{belkin2019reconciling,nakkiran2019double} 顯示，模型大到超過插值門檻之後，再變大反而泛化更好。頓悟 \cite{power2022grokking} 顯示，小模型在長時間訓練後可能從記憶轉為學會規則；Wang 等人 \cite{wang2024grokked} 進一步發現兩步推理能否學會，取決於推理例子與基本事實的比例。捷徑學習 \cite{geirhos2020shortcut} 說明模型傾向抓取最容易的線索。組合泛化的研究（SCAN \cite{lake2018generalization}、COGS \cite{kim2020cogs}）顯示 Transformer 不擅長新的組合。

\paragraph{知識的儲存與取用。}
Allen-Zhu 與 Li \cite{allenzhu2023physics31} 發現，同一事實只用一種寫法出現時，模型很難在問答中取用；用多種寫法並打亂順序後，對沒練過問答的人物也能答對大部分。他們也估計每個參數大約能儲存 2 bits 的知識 \cite{allenzhu2024physics33}。反轉詛咒 \cite{berglund2023reversal} 顯示大模型學了「A 是 B」，常答不出反方向的問題。

\paragraph{湧現式溝通。}
Lazaridou 等人 \cite{lazaridou2017multi}、Havrylov 與 Titov \cite{havrylov2017emergence}、Mordatch 與 Abbeel \cite{mordatch2018emergence} 讓多個代理人透過指稱遊戲發展出自己的溝通方式。本文的 mini-talk 不發明新語言，而是讓兩個模型用同一種受限語言交換知識。

% ============================================================
\section{Mini-Lang 語言設計}\label{sec:minilang}

\subsection{設計原則}

表 \ref{tab:principles} 列出 Mini-Lang 的設計原則。其中最重要的是「子集」性質：每一句合乎規則的 Mini-Lang 句子，都是母語者聽得懂、而且自然的標準中文。因此學會 Mini-Lang 的人（或模型）說出的話，可以直接和一般中文使用者溝通。

\begin{table}[H]
\centering
\caption{Mini-Lang 的設計原則}\label{tab:principles}
\begin{tabularx}{\linewidth}{lY}
\toprule
原則 & 做法 \\
\midrule
字少 & 全部 256 字，即 $2^8$，可以用一個位元組編碼 \\
是真正的中文 & 每一句合乎規則的句子，都是自然的標準中文 \\
詞靠組合 & 字是積木，詞由中文本來就有的構詞方式組成（電＋腦＝電腦） \\
語法固定 & 一條固定語序與 17 條句型規則，不用「把」字句、「被」字句等結構 \\
換句話說 & 字表外的詞改寫，例如「便宜 $\to$ 不貴」「醫院 $\to$ 看病的地方」 \\
\bottomrule
\end{tabularx}
\end{table}

\subsection{字表}

256 字分成九類（表 \ref{tab:charset}），完整字表見附錄 \ref{app:charset}。

\begin{table}[H]
\centering
\caption{Mini-Lang 字表的九個類別}\label{tab:charset}
\begin{tabularx}{\linewidth}{lrY}
\toprule
類別 & 字數 & 例子 \\
\midrule
代詞與指示 & 15 & 我、你、他、這、那、什麼、怎麼樣 \\
語法 & 41 & 的、了、過、嗎、不、沒、是、有、在、因為、所以、如果 \\
數量 & 23 & 一到十、百、千、萬、兩、幾、個、些、次、點、塊、歲 \\
時間 & 17 & 年、月、日、天、今、明、昨、早、晚、分、鐘、週 \\
方位 & 17 & 上下左右前後、裡外中間、東南西北、邊、方、地 \\
動作 & 50 & 來、去、吃、喝、說、聽、看、寫、學、買、賣、找、幫 \\
性質 & 30 & 大、小、多、少、好、壞、新、快、慢、冷、熱、難、容易 \\
事物 & 59 & 人、家人、身體、水、火、山、車、書、錢、飯、茶、國、語言 \\
社交 & 4 & 請、謝、迎、係 \\
\midrule
合計 & 256 & \\
\bottomrule
\end{tabularx}
\end{table}

\subsection{語法}

句子依照圖 \ref{fig:slots} 的固定位置排列，括號內的位置可以省略。17 條規則涵蓋主謂賓、時間與地點、是字句、有字句、性質謂語、修飾、數量、否定、問句、時態、能願動詞、比較、方向、得字補語、介詞片語、連接詞與祈使句。

\begin{figure}[H]
\centering
\begin{tikzpicture}[
  slot/.style={draw, rounded corners=2pt, minimum height=6mm, inner sep=2.5pt, font=\footnotesize},
  opt/.style={slot, dashed},
  arr/.style={-{Stealth[length=2mm]}},
  node distance=2mm]
\node[opt] (t) {時間};
\node[slot, right=of t] (s) {主語};
\node[opt, right=of s] (ad) {也／都／還／就};
\node[opt, right=of ad] (n) {不／沒};
\node[opt, right=of n] (m) {能願};
\node[opt, below=8mm of t, xshift=4mm] (p) {介詞片語};
\node[slot, right=of p] (v) {動詞};
\node[opt, right=of v] (c) {補語／了／過};
\node[opt, right=of c] (o) {賓語};
\node[opt, right=of o] (e) {語氣詞};
\draw[arr] (t) -- (s); \draw[arr] (s) -- (ad); \draw[arr] (ad) -- (n); \draw[arr] (n) -- (m);
\draw[arr] (m.south) |- ++(0,-4mm) -| (p.north);
\draw[arr] (p) -- (v); \draw[arr] (v) -- (c); \draw[arr] (c) -- (o); \draw[arr] (o) -- (e);
\end{tikzpicture}
\caption{Mini-Lang 句子的固定位置。實線框為必要位置，虛線框可以省略。}\label{fig:slots}
\end{figure}

為了維持「子集」的性質，我們做了幾個取捨：第三人稱統一用「他」；只用個、些、次、點、塊、歲六個量詞，需要特殊量詞的名詞不直接加數字（說「一些書」而不是「一本書」）；用「週一」代替「星期一」。

\subsection{自動檢查}

\texttt{check.py} 讀取規格文件中的字表，確認字表正好 256 字且不重複，並檢查文件中所有例句只用到字表內的字。之後所有語料都經過同一個工具檢查。

% ============================================================
\section{系統}\label{sec:system}

\subsection{模型}

我們使用類似 LLaMA 的字元級 Transformer：RMSNorm、旋轉位置編碼（RoPE）、SwiGLU 前饋層，輸入嵌入與輸出層共用權重。詞表由所有字元加上五個特殊 token 組成：\tok{PAD}、\tok{UNK}、\tok{EOS}、\tok{Q}、\tok{A}。表 \ref{tab:models} 列出實驗中使用的四種模型大小。

\begin{table}[H]
\centering
\caption{模型大小（參數量以 Mini-Lang 的 269 個 token 詞表計算）}\label{tab:models}
\begin{tabular}{lrrrr}
\toprule
名稱 & $d_{\text{model}}$ & 層數 & 頭數 & 參數量 \\
\midrule
tiny & 32 & 2 & 2 & 41,536 \\
small & 64 & 2 & 4 & 148,608 \\
base & 128 & 4 & 4 & 1,084,160 \\
large & 256 & 6 & 8 & 6,363,648 \\
\bottomrule
\end{tabular}
\end{table}

\subsection{訓練}

訓練分兩階段。\textbf{預訓練}時，每句結尾加上 \tok{EOS}，每個樣本都重新打亂句子順序再拼接，以免模型記住句子之間的固定順序。\textbf{微調}時，序列為「\tok{Q}問題\tok{A}答案\tok{EOS}」，只對「答案＋\tok{EOS}」計算 loss；在 mini-talk 中，問題部分也計算 loss，模型才能學會發問。

生成時遇到 \tok{EOS} 即停止，評估一律用 greedy 解碼，答案完全一致才算對。除非另外說明，預訓練 2,000 步、微調 500 步、batch 32、序列長度 96，使用 AdamW（預訓練學習率 $5\times10^{-4}$、微調 $1\times10^{-4}$，weight decay 0.01），全部在 Apple M3 的 CPU 上執行。

\subsection{語料}

表 \ref{tab:corpora} 列出本文使用的四種語料。

\begin{table}[H]
\centering
\caption{語料}\label{tab:corpora}
\begin{tabularx}{\linewidth}{lYYl}
\toprule
語料 & 來源 & 規模 & 用途 \\
\midrule
planets & 大型語言模型蒸餾的八大行星知識 & 50 句事實、50 組問答 & 起點 \\
minilang & 程式產生：人物事實、語法規則句、文件例句 & 16 人時 1,802 句、112 組問答 & 主要實驗 \\
handwritten & 刻意使用模板沒有的句型（由 Claude 撰寫） & 118 句訓練、40 句測試 & 自然句子 \\
directions & 大型語言模型依 13 個方向撰寫，經字表過濾 & 1,805 句訓練、196 句測試 & 覆蓋各方向 \\
\bottomrule
\end{tabularx}
\end{table}

\paragraph{人物世界。}
產生器建立 $N$ 個人物，每人有 7 項屬性：年齡、工作、住處、愛吃、愛喝、朋友、孩子數。名字也由字表內的字組成（如「小明」「大山」）。預訓練語料包含這些事實的敘述句，微調語料包含對應的問答。驗證時，保留 20\% 的問答不微調，但它們的事實在預訓練中出現過，用來測試模型能否把預訓練的知識用在問答上。

\paragraph{方向語料。}
我們讓七個 Claude 子代理分別撰寫 13 個方向的語句：日常社交、多輪問答、時間與數量、買東西與吃喝、地點與方向、交通、家庭人物與工作、健康與身體、學習與語言、天氣與自然、情感與看法、推理與因果、短篇故事。原始 2,015 行經過字表檢查、人名過濾（排除 6 行）與去重（排除 8 行）後，每個方向保留約 10\% 作為測試。

\subsection{mini-talk}

為了研究兩個模型能否透過對話交換知識，我們把 16 個人物分成兩組。模型 A 只在預訓練中看過 A 組 8 人的事實，模型 B 只看過 B 組的事實；兩者都學了相同的一般句子。微調時兩者都看全部 16 人的問題，自己知道的人回答事實，不知道的人回答「我不知道」。交談規則如下：
\begin{enumerate}
  \item 發問者從 \tok{Q} 開始抽樣產生問題，生成到 \tok{A} 為止。
  \item 發問者先問自己；如果自己答得出來，就重新產生問題，也就是只問自己不知道的事。
  \item 對方用 greedy 解碼回答。
  \item 交談後，每個模型用聽到的答案再微調，並混入原本的問答以避免遺忘。
\end{enumerate}

% ============================================================
\section{實驗}\label{sec:exp}

\subsection{起點：小資料的死背}

在 planets 語料上（約 108 萬參數），訓練集 40/40 全對，驗證集 0/10。模型能背下所有問答，卻完全無法把事實敘述轉為問答。

改用 minilang 語料（16 人、base 模型）後，訓練集 90/90，驗證集 11/22（50\%）。同一種「事實 $\to$ 問答」的對應有 16 個人可以練習，模型開始能推廣到微調沒看過的問題。

\subsection{縮放實驗：模型大小 $\times$ 資料量}

我們把人物數設為 8、16、64、256（0.5、1、4、16 倍），一般句子相應增加。模板最多只能產生約 9,600 種不同的一般句子，因此 4 倍與 16 倍時句子數分別為 5,180 與 8,441。基準是「每種問題都回答訓練集中最常見的答案」。結果見表 \ref{tab:scaling} 與圖 \ref{fig:scaling}。

\begin{table}[H]
\centering
\caption{驗證集答對率（新人物事實）。正則化組為 dropout 0.1、weight decay 0.1。}\label{tab:scaling}
\small
\begin{tabular}{lrrrrr}
\toprule
模型 & 參數 & 0.5 倍 & 1 倍 & 4 倍 & 16 倍 \\
\midrule
tiny & 41,536 & 3/11 (27.3\%) & 5/22 (22.7\%) & 20/89 (22.5\%) & 68/358 (19.0\%) \\
small & 148,608 & 5/11 (45.5\%) & 7/22 (31.8\%) & 32/89 (36.0\%) & 54/358 (15.1\%) \\
base & 1,084,160 & 6/11 (54.5\%) & 11/22 (50.0\%) & 71/89 (79.8\%) & 325/358 (90.8\%) \\
large & 6,363,648 & 7/11 (63.6\%) & 17/22 (77.3\%) & 81/89 (91.0\%) & \textbf{353/358 (98.6\%)} \\
\midrule
base＋正則化 & 1,084,160 & 9/11 (81.8\%) & 14/22 (63.6\%) & 70/89 (78.7\%) & 121/358 (33.8\%) \\
最常見答案 & --- & 3/11 (27\%) & 7/22 (32\%) & 17/89 (19\%) & 88/358 (25\%) \\
\bottomrule
\end{tabular}
\end{table}

\begin{figure}[H]
\centering
\begin{tikzpicture}[x=2.6cm, y=0.045cm, font=\small]
  % 座標軸
  \draw[->] (0,0) -- (3.4,0) node[right] {資料倍數};
  \draw[->] (0,0) -- (0,108) node[above] {答對率（\%）};
  \foreach \y in {0,20,40,60,80,100} {
    \draw (0,\y) -- (-0.04,\y) node[left] {\y};
    \draw[gray!25] (0,\y) -- (3.2,\y);
  }
  \foreach \x/\l in {0.2/0.5, 1.2/1, 2.2/4, 3.2/16} \draw (\x,0) -- (\x,-2) node[below] {\l};
  % 各模型
  \draw[thick, gray]  plot coordinates {(0.2,27.3) (1.2,22.7) (2.2,22.5) (3.2,19.0)};
  \draw[thick, orange] plot coordinates {(0.2,45.5) (1.2,31.8) (2.2,36.0) (3.2,15.1)};
  \draw[thick, blue]  plot coordinates {(0.2,54.5) (1.2,50.0) (2.2,79.8) (3.2,90.8)};
  \draw[thick, red]   plot coordinates {(0.2,63.6) (1.2,77.3) (2.2,91.0) (3.2,98.6)};
  \draw[thick, blue, dashed] plot coordinates {(0.2,81.8) (1.2,63.6) (2.2,78.7) (3.2,33.8)};
  \draw[black, dotted, thick] plot coordinates {(0.2,27) (1.2,32) (2.2,19) (3.2,25)};
  \foreach \p in {(0.2,63.6),(1.2,77.3),(2.2,91.0),(3.2,98.6)} \fill[red] \p circle (1.5pt);
  \foreach \p in {(0.2,54.5),(1.2,50.0),(2.2,79.8),(3.2,90.8)} \fill[blue] \p circle (1.5pt);
  % 圖例
  \begin{scope}[shift={(3.45,95)}, font=\footnotesize]
    \draw[thick, red] (0,0) -- (0.2,0) node[right, black] {large};
    \draw[thick, blue] (0,-10) -- (0.2,-10) node[right, black] {base};
    \draw[thick, orange] (0,-20) -- (0.2,-20) node[right, black] {small};
    \draw[thick, gray] (0,-30) -- (0.2,-30) node[right, black] {tiny};
    \draw[thick, blue, dashed] (0,-40) -- (0.2,-40) node[right, black] {base＋正則化};
    \draw[thick, dotted] (0,-50) -- (0.2,-50) node[right, black] {最常見答案};
  \end{scope}
\end{tikzpicture}
\caption{縮放實驗：驗證集答對率隨資料量的變化（橫軸不等距）。}\label{fig:scaling}
\end{figure}

我們得到四項發現：
\begin{enumerate}
  \item \textbf{縮小模型在每一種資料量下都更差。}每一欄都是 large $>$ base $>$ small $>$ tiny。我們沒有觀察到雙重下降的高峰，可能是模型大小的間隔太大。
  \item \textbf{最小的模型只是在猜。}tiny 的成績和最常見答案基準相當：它學會了答案的格式，卻沒學會取用知識。
  \item \textbf{增加資料要搭配足夠的容量。}base 與 large 隨資料增加大幅進步；tiny 與 small 卻隨資料增加而變差，small 在 16 倍時甚至低於基準。在本實驗中，資料增加也代表要記的人物增加，小模型的容量不足以應付。
  \item \textbf{正則化只在資料少時有益。}0.5 倍時提升 27 個百分點，16 倍時下降 57 個百分點。資料少時模型傾向死背，正則化把它推向規則；資料多時模型需要全部的容量，正則化造成欠擬合。
\end{enumerate}

\subsection{資料設計：定向覆蓋}

我們固定使用 large 模型與 256 人的世界，比較表 \ref{tab:configs} 的四種語料設定。A 即縮放實驗中的 large $\times$ 16 倍（微調 500 步）；B、C、D 微調 1,000 步。每種問題準備 3 種問法，其中第 3 種完全不用於微調。

\begin{table}[H]
\centering
\caption{資料設計實驗的四種設定}\label{tab:configs}
\begin{tabularx}{\linewidth}{lYl}
\toprule
設定 & 人物資料 & 其他語句 \\
\midrule
A 基準 & 每個事實 2 種說法，每種問題 1 種問法 & 模板句子 \\
B 加強資料 & 每個事實 4 種說法，每人 2 段打亂順序的小傳，每種問題 2 種問法 & 模板句子 \\
C 加強＋手寫 & 同 B & ＋handwritten \\
D 加強＋手寫＋方向 & 同 B & ＋handwritten＋directions \\
\bottomrule
\end{tabularx}
\end{table}

我們用四種測試評估泛化（表 \ref{tab:tests}）。均勻亂猜的每字交叉熵為 $\ln 269 \approx 5.59$。

\begin{table}[H]
\centering
\caption{四種泛化測試}\label{tab:tests}
\begin{tabularx}{\linewidth}{lYr}
\toprule
測試 & 內容 & 題數 \\
\midrule
新事實 & 沒微調過的事實，原本的問法（與表 \ref{tab:scaling} 的驗證集相同） & 358 \\
新問法 & 微調過的事實，從未出現的問法（如「誰是小明的朋友？」） & 1,434 \\
兩者皆新 & 沒微調過的事實＋從未出現的問法 & 358 \\
新句子 & 手寫測試句與方向測試句的每字交叉熵 & 40＋196 句 \\
\bottomrule
\end{tabularx}
\end{table}

\begin{table}[H]
\centering
\caption{資料設計實驗結果。方向句 loss 為 13 個方向的平均。}\label{tab:best}
\begin{tabular}{lrrrrr}
\toprule
設定 & 新事實 & 新問法 & 兩者皆新 & 手寫句 loss & 方向句 loss \\
\midrule
A 基準 & \textbf{98.6\%} & 68.5\% & 57.5\% & 5.54 & 6.74 \\
B 加強資料 & 80.7\% & \textbf{85.5\%} & \textbf{64.2\%} & 5.91 & 6.92 \\
C 加強＋手寫 & 83.5\% & 73.7\% & 53.1\% & 4.89 & 5.99 \\
D 加強＋手寫＋方向 & 79.1\% & 80.6\% & 52.0\% & \textbf{2.35} & \textbf{3.30} \\
\bottomrule
\end{tabular}
\end{table}

D 在每一個方向都把 loss 降到 A 的一半左右（表 \ref{tab:directions}）。表 \ref{tab:samples} 用相同開頭讓預訓練模型續寫：A 只會模板的句子，前後接不起來；D 能接出合理的回應，甚至寫出多句的連貫內容。

\begin{table}[H]
\centering
\caption{各方向測試句的每字 loss}\label{tab:directions}
\small
\begin{tabular}{lrrrr}
\toprule
方向 & A & B & C & D \\
\midrule
日常社交 & 6.58 & 6.74 & 5.42 & \textbf{3.17} \\
多輪問答 & 7.64 & 8.52 & 6.87 & \textbf{3.42} \\
時間與數量 & 7.30 & 7.45 & 6.08 & \textbf{3.76} \\
買東西與吃喝 & 6.22 & 6.26 & 5.97 & \textbf{3.45} \\
地點與方向 & 6.84 & 7.50 & 6.24 & \textbf{3.19} \\
交通 & 7.14 & 7.09 & 6.40 & \textbf{3.48} \\
家庭人物與工作 & 6.05 & 6.25 & 5.38 & \textbf{2.61} \\
健康與身體 & 6.61 & 6.52 & 5.61 & \textbf{3.25} \\
學習與語言 & 6.63 & 6.86 & 6.16 & \textbf{3.33} \\
天氣與自然 & 6.45 & 6.41 & 5.86 & \textbf{3.23} \\
情感與看法 & 6.40 & 6.41 & 5.14 & \textbf{3.02} \\
推理與因果 & 6.12 & 6.32 & 5.86 & \textbf{3.41} \\
短篇故事 & 7.67 & 7.62 & 6.81 & \textbf{3.54} \\
\bottomrule
\end{tabular}
\end{table}

\begin{table}[H]
\centering
\caption{用相同開頭讓預訓練模型續寫（溫度 0.7，top-$k$ 10）}\label{tab:samples}
\begin{tabularx}{\linewidth}{lYY}
\toprule
開頭 & A 基準 & D \\
\midrule
今天早上下雨了， & 我姐姐就在家裡看書。 & 我們回家吧。 \\
請問，看病的地方 & 山沒有孩子。 & 不在這邊，在那邊。 \\
我的手機壞了， & 我在學校學習吧。 & 我不知道了。可是那裡的人都很好，有一個人還幫我們找到了住的地方。 \\
\bottomrule
\end{tabularx}
\end{table}

\paragraph{新事實準確率下降的分析。}
B 在「新事實」從 98.6\% 降到 80.7\%，而且七類問題全部退步。為了找出原因，我們做了兩個檢查。第一，把問題改寫成事實句的開頭（如「小雨喜歡吃」）讓預訓練模型續寫：A 正確 270/358（75\%），B 正確 259/358（72\%），兩者記得的事實相近。第二，把 B 的微調從 1,000 步減為 500 步：新事實回升到 306/358（85\%），新問法降到 1,167/1,434（81\%），兩者皆新為 210/358（59\%）。因此，微調過久只解釋了一小部分差距。我們推測主因是相同的訓練量被分散到更多說法與問法上，使「原本問法」這個特定方向的表現下降。

\subsection{對話學習：mini-talk}

兩個 base 模型各自預訓練 2,000 步、微調 500 步（問題部分也計算 loss），每人各問 20 題，結果見表 \ref{tab:talk}。「學習後」是對方那組人的 56 題中答對的題數，學習前兩邊都是 0/56。

\begin{table}[H]
\centering
\caption{mini-talk 結果}\label{tab:talk}
\begin{tabular}{rrrrrrr}
\toprule
發問溫度 & 句子數 & 字表外 & 不通順的問題 & 回答正確 & A 學習後 & B 學習後 \\
\midrule
1.0 & 80 & 0 & 5 & 35/35 & 17/56 & 18/56 \\
0.7 & 80 & 0 & 0 & 40/40 & 20/56 & 20/56 \\
0.5 & 80 & 0 & 0 & 40/40 & 20/56 & 20/56 \\
\bottomrule
\end{tabular}
\end{table}

所有句子都只用到 256 字；問題通順時，回答全部正確。學習後答對的題數正好等於聽過的題數（20/20），沒聽過的同一人物的其他事實仍然答不出來。對話確實傳遞了資訊，但模型記住的是「問題 $\to$ 答案」，而不是「關於這個人的知識」。溫度 1.0 時，模型產生了不通順的問題，對方卻照樣作答，例如「小明的外在哪裡？」被回答成「南邊」。這表示模型依據局部線索（人名、「哪裡」）挑選答案，而非理解整句話。

% ============================================================
\section{討論}\label{sec:discussion}

\subsection{泛化取決於容量、資料量與資料變化的關係}

縮放實驗顯示，「小模型泛化差」只說對一半。原本的設定（108 萬參數、約 2 萬字資料）中，資料量只有參數量的約 0.02 倍，模型能把所有資料背下來，死背比學規則便宜。增加資料的量與變化讓死背不再划算，模型才被迫學規則；但這需要足夠的容量，資料增加時，容量不足的模型反而變差。這和雙重下降的觀點一致：在已經過度參數化的區域，模型越大，泛化越好。也和 Kaplan 等人 \cite{kaplan2020scaling} 對樣本效率的觀察一致：資料有限時，大模型從每筆資料學到更多。

\subsection{定向覆蓋}

資料設計實驗顯示，加入的語句帶來哪一種變化，模型就在那個方向學會推廣：同義問法讓模型在換一種問法時答得出來；手寫生活句讓沒看過的自然句子的 loss 下降；13 個方向的語句讓所有方向的新句子 loss 減半。反過來，只用模板訓練的 A、B，對自然句子的 loss（5.5--6.9）甚至高於均勻亂猜（5.59），表示它們不只是不會，而是很有把握地猜錯。

如果方向有 $k$ 個、每個方向有 $n$ 種變化，各自覆蓋只需約 $nk$ 個例子，所有組合卻有 $n^k$ 種。因此可行的策略不是列舉組合，而是讓每個方向各自有足夠的變化，再測試模型能否組合。「兩者皆新」測試就是在檢驗這一點：B 的表現（64.2\%）比 A（57.5\%）好，但仍明顯低於各自單一方向的成績。

Mini-Lang 的價值在於方向是數得完的：256 字、17 條規則、有限的生活情境。在一般中文上很難做到的完整覆蓋，在 Mini-Lang 上可以逐一檢查。

\subsection{取捨}

加入更多形式提升了換句話問的能力，卻降低了「原本問法＋新事實」的準確率。在固定訓練量下，資料多樣性與單一方向的準確率之間存在取捨。可能的緩解方式包括增加訓練步數、微調時混入預訓練語料（replay），以及進一步加大模型。

\subsection{理解}

mini-talk 中「不通順的問題照樣作答」，以及「只記得聽過的題目」，都顯示模型的行為更接近依局部線索查表，而非理解。我們認為理解只能從行為推斷，因此需要一組由易到難的測試：換句話問、換人問、問不存在的事、兩步推理、置換答案的因果測試，以及以 slot filling（QA-SRL \cite{he2015qasrl} 形式）明確寫出模型對句子的解讀。本文只完成了其中「換句話問」的部分。

% ============================================================
\section{限制}

\begin{enumerate}
  \item \textbf{單一亂數種子。}每個設定只訓練一次。B 與 C 只差不到 1\% 的手寫句子，新問法卻差了 11.8 個百分點，表示訓練本身的隨機誤差很大。本文只採用差距大、方向一致的結論。
  \item \textbf{計算量有限。}所有實驗都在 CPU 上執行，預訓練固定 2,000 步，語料較大的設定可能沒有訓練充分。
  \item \textbf{測試資料與訓練資料的來源相近。}方向語料的測試句與訓練句由同一批子代理撰寫，風格相近，因此 D 在方向句上的進步有一部分是分布內的泛化。手寫測試句雖然獨立撰寫，作者同樣是 Claude。
  \item \textbf{人物世界過於簡單。}只有 7 種屬性，每種屬性的值很少，答案種類有限，亂猜的基準也因此偏高（19--32\%）。
  \item \textbf{Mini-Lang 的取捨。}統一用「他」、只用六個量詞等規則，使部分句子和一般中文的習慣略有不同。
\end{enumerate}

% ============================================================
\section{結論與未來工作}

本文提出 Mini-Lang，一個 256 字、規則固定、且是自然中文子集的受限語言，並以它為平台研究小型語言模型的泛化。實驗顯示：縮小模型不會讓它「被迫學規則」，反而在每一種資料量下都更差；增加資料必須搭配足夠的容量；正則化只在資料少時有益；混入的語句只幫助它所代表的方向，而以大型語言模型依方向產生的少量語料，就能大幅提升模型對沒看過的句子的處理能力。

未來工作包括：以多個亂數種子重複實驗以估計誤差；在微調中加入 replay 以減少知識遺忘；擴充方向語料至每個方向數千句，並加入拒答與兩步推理問答；建立 slot filling 與其他理解測試；讓 mini-talk 中的模型根據前文追問，並以 slot 一致率測量溝通成功率。

\section*{致謝與聲明}

本研究的程式、Mini-Lang 規格文件、手寫語料與方向語料，均在 Claude（Anthropic）的協助下完成。方向語料由 Claude 的子代理依方向撰寫，再經 \texttt{check.py} 與 \texttt{build\_directions.py} 自動檢查。所有實驗數據都來自實際執行的結果，原始數據保存於專案中的 \texttt{scaling\_results.csv} 與 \texttt{best\_model\_results.csv}。

\section*{可重現性}

\begin{table}[H]
\centering
\small
\begin{tabularx}{\linewidth}{lY}
\toprule
項目 & 位置 \\
\midrule
Mini-Lang 規格與檢查工具 & \texttt{mini-lang/mini-lang.md}、\texttt{mini-lang/check.py} \\
模型與訓練程式 & \texttt{mini-lm/model.py}、\texttt{pretrain.py}、\texttt{finetune.py}、\texttt{trainer.py} \\
語料產生 & \texttt{mini-lm/gen\_minilang.py}、\texttt{build\_directions.py}、\texttt{corpus/} \\
縮放實驗 & \texttt{mini-lm/experiments/scaling.py} \\
資料設計實驗 & \texttt{mini-lm/experiments/best\_model.py} \\
對話實驗 & \texttt{mini-talk/} \\
\bottomrule
\end{tabularx}
\end{table}

% ============================================================
\begin{thebibliography}{99}
\small
\bibitem{allenzhu2023physics31} Z. Allen-Zhu and Y. Li. Physics of Language Models: Part 3.1, Knowledge Storage and Extraction. \textit{arXiv:2309.14316}, 2023.
\bibitem{allenzhu2024physics33} Z. Allen-Zhu and Y. Li. Physics of Language Models: Part 3.3, Knowledge Capacity Scaling Laws. \textit{arXiv:2404.05405}, 2024.
\bibitem{belkin2019reconciling} M. Belkin, D. Hsu, S. Ma, and S. Mandal. Reconciling modern machine-learning practice and the classical bias--variance trade-off. \textit{PNAS}, 116(32), 2019.
\bibitem{berglund2023reversal} L. Berglund et al. The Reversal Curse: LLMs trained on ``A is B'' fail to learn ``B is A''. \textit{arXiv:2309.12288}, 2023.
\bibitem{eldan2023tinystories} R. Eldan and Y. Li. TinyStories: How Small Can Language Models Be and Still Speak Coherent English? \textit{arXiv:2305.07759}, 2023.
\bibitem{geirhos2020shortcut} R. Geirhos et al. Shortcut Learning in Deep Neural Networks. \textit{Nature Machine Intelligence}, 2, 2020.
\bibitem{havrylov2017emergence} S. Havrylov and I. Titov. Emergence of Language with Multi-agent Games: Learning to Communicate with Sequences of Symbols. \textit{NeurIPS}, 2017.
\bibitem{he2015qasrl} L. He, M. Lewis, and L. Zettlemoyer. Question-Answer Driven Semantic Role Labeling. \textit{EMNLP}, 2015.
\bibitem{hoffmann2022chinchilla} J. Hoffmann et al. Training Compute-Optimal Large Language Models. \textit{arXiv:2203.15556}, 2022.
\bibitem{kaplan2020scaling} J. Kaplan et al. Scaling Laws for Neural Language Models. \textit{arXiv:2001.08361}, 2020.
\bibitem{kim2020cogs} N. Kim and T. Linzen. COGS: A Compositional Generalization Challenge Based on Semantic Interpretation. \textit{EMNLP}, 2020.
\bibitem{lake2018generalization} B. Lake and M. Baroni. Generalization without Systematicity. \textit{ICML}, 2018.
\bibitem{lazaridou2017multi} A. Lazaridou, A. Peysakhovich, and M. Baroni. Multi-Agent Cooperation and the Emergence of (Natural) Language. \textit{ICLR}, 2017.
\bibitem{mordatch2018emergence} I. Mordatch and P. Abbeel. Emergence of Grounded Compositional Language in Multi-Agent Populations. \textit{AAAI}, 2018.
\bibitem{nakkiran2019double} P. Nakkiran et al. Deep Double Descent: Where Bigger Models and More Data Hurt. \textit{arXiv:1912.02292}, 2019.
\bibitem{ogden1930basic} C. K. Ogden. \textit{Basic English: A General Introduction with Rules and Grammar}. Kegan Paul, 1930.
\bibitem{power2022grokking} A. Power et al. Grokking: Generalization Beyond Overfitting on Small Algorithmic Datasets. \textit{arXiv:2201.02177}, 2022.
\bibitem{schaeffer2023mirage} R. Schaeffer, B. Miranda, and S. Koyejo. Are Emergent Abilities of Large Language Models a Mirage? \textit{NeurIPS}, 2023.
\bibitem{wang2024grokked} B. Wang et al. Grokked Transformers are Implicit Reasoners: A Mechanistic Journey to the Edge of Generalization. \textit{arXiv:2405.15071}, 2024.
\bibitem{wei2022emergent} J. Wei et al. Emergent Abilities of Large Language Models. \textit{TMLR}, 2022.
\end{thebibliography}

% ============================================================
\appendix
\section{Mini-Lang 完整字表}\label{app:charset}

\begin{tabularx}{\linewidth}{lY}
\toprule
類別 & 字 \\
\midrule
代詞與指示 & 我你他們這那誰什麼哪怎樣自己每 \\
語法 & 的了過嗎呢吧不沒是有在和也都很太最更還就再會能要可以得給對跟從到為因所如果但只起比 \\
數量 & 一二三四五六七八九十百千萬兩幾半第個些次點塊歲 \\
時間 & 年月日天時候今明昨早晚午現分鐘週先 \\
方位 & 上下左右前後裡外中間東西南北邊方地 \\
動作 & 來去走坐站住回進出開關拿放用做作說話聽看見寫學習教問叫想知道懂記覺喜歡愛吃喝睡買賣找等幫穿洗玩打生完 \\
性質 & 大小多少長短高老好壞新快慢冷熱難容易真忙病重遠近白黑紅色樂貴 \\
事物 & 人男女子孩爸媽哥姐弟妹朋友師家名字身體頭手眼口心水火山雨氣電車路門房書衣服錢飯菜茶肉花校醫店國文語言事情題工腦機飛意思 \\
社交 & 請謝迎係 \\
\bottomrule
\end{tabularx}

\section{方向語料範例}

\begin{tabularx}{\linewidth}{lY}
\toprule
方向 & 範例 \\
\midrule
多輪問答 & 「你有沒有時間？」「有，什麼事？」「你可以幫我一個忙嗎？」「可以，你說吧。」 \\
推理與因果 & 如果他是醫生，他就會看病，他不會看病，所以他不是醫生。 \\
家庭人物與工作 & 我哥哥是我爸爸媽媽的孩子，我也是，所以我們是一家人。 \\
學習與語言 & 「電」和「腦」放在一起，就是「電腦」。 \\
短篇故事 & 我的手機不見了。我在家裡找了一個小時，都沒有找到。後來我想到，我可以用家裡的電話打給我的手機。電話一打，手機就在衣服裡邊叫起來了。 \\
\bottomrule
\end{tabularx}

\end{document}
